\begin{definition}{Banach--Mazur--Abstand, Banach--Mazur--Abbildung}
                  {BanachMazurAbstand}
Seien $(E, \norm\cdot_E), (F, \norm\cdot_F)$ normierte $\K$--Vektorr�ume.\\
Dann hei�t
\[d_{BM}\left( (E, \norm\cdot_E), (F, \norm\cdot_F) \right) := 
\inf\menge{\opnorm S E F \cdot \opnorm {S\inv} F E}{S \in \Inv(E, F)}\]
\textbf{Banach--Mazur--Abstand zwischen $(E, \norm\cdot_E)$ und 
$(F, \norm\cdot_F)$}.\\
Wenn aus dem Zusammenhang klar ist, welche Norm auf $E$ bzw. $F$ vorliegt,
werden wir auch $d_{BM}(E,F)$ schreiben.\\
Durch oben erw�hnte Festsetzung $\inf \emptyset = \infty$ ist $d_{BM}$ f�r je
zwei normierte R�ume sinnvoll erkl�rt.\\
Ist $\Inv(E, F) \neq \emptyset$, so hei�t ein $T \in \Inv(E, F)$ \textbf{Banach--Mazur--Abbildung zwischen $(E, \norm\cdot_E)$ und 
$(F, \norm\cdot_F)$}
\[:\iff \opnorm T E F = 1 \quad \wedge \quad 
\opnorm {T\inv} F E 
 = d_{BM}\left( (E, \norm\cdot_E), (F, \norm\cdot_F) \right)\text.\]
Insbesondere gilt also f�r eine Banach--Mazur--Abbildung $T$:
\[\opnorm T E F \cdot \opnorm {T\inv} F E 
= d_{BM}\left( (E, \norm\cdot_E), (F, \norm\cdot_F) \right)\text.\]
\end{definition}